\begin{abstract}
We introduce \method, a graph prompting method that uses support-set prototypes to route trainable prompts to individual nodes. With only a few labeled nodes per class, adaptation must determine both what a prompt should encode and where its contribution should be applied. \method{} constructs class references from aggregated node attributes, refines them using unlabeled graph features while retaining support anchors, and converts node--prototype associations into a fixed routing matrix. The resulting combinations of class prompt tokens are injected into a frozen graph convolutional encoder, followed by prototype classification. We evaluate each 1-, 3-, and 5-shot setting separately on Cora, CiteSeer, and PubMed using matched support draws and training seeds. In the 5-shot setting, \method{} achieves accuracies of 73.02\%, 66.17\%, and 70.98\%, respectively; replacing a shared prompt with prototype routing improves accuracy by 36.94, 40.45, and 3.02 percentage points in the corresponding conditions. These paired results support the value of prototype-guided prompt allocation within the evaluated transductive setting.
\end{abstract}

\section{Introduction}
Pretrained graph neural networks offer a starting point for learning with scarce node labels. Graph prompting adapts these representations through a small set of task parameters, using feature modifications, prompt structures, or task-specific readout~\citep{gppt2022,gpf2023,allinone2023,graphprompt2023}. Beyond choosing an injection point, an adaptation method must decide how a prompt's contribution is distributed across nodes. This allocation matters because nodes in the same graph can play different roles in the downstream classification task.

A shared prompt expresses a global adjustment. More flexible methods learn input-dependent mixtures, assign prompts to subgraphs, or select nodes through a policy~\citep{gpf2023,supt2026,relief2025}. These approaches establish that node-dependent allocation is not itself a new idea. Our question is more specific: \emph{can the few labeled examples already available for a task define useful routing coefficients, so that optimization can focus on prompt content?} Support examples provide class references that can be constructed before adaptation and used without training an additional routing network.

We propose \method{}, which uses support-anchored prototypes as coordinates for prompt allocation. Nodes with similar class associations receive related combinations of class prompt tokens. The routing representation comes from aggregated graph attributes; its prototypes are initialized by labeled support nodes and refined with unlabeled features while retaining the support anchors. These associations yield a fixed routing matrix that mixes trainable tokens inside a frozen GCN. Only the prompt tokens are optimized, and prediction uses a prototype readout. Figure~\ref{fig:architecture} summarizes this separation between fixed allocation and learned content.

Our evaluation asks two questions. First, how does complete adaptation compare with frozen pretraining and full fine-tuning in matched local runs, and where does its accuracy lie relative to literature-reported prompt results? Second, does prototype routing add value over a shared prompt when the injection point and training recipe are held fixed? Local runs cover Cora, CiteSeer, and PubMed at 1, 3, and 5 shots. External accuracy values come from \citet{mtg2026} and retain their original protocol. Every result remains specific to its dataset and shot count. The gain over frozen pretraining is much smaller on PubMed than on Cora or CiteSeer.

The paper makes three contributions:
\begin{itemize}
\item A graph prompting construction that uses support-anchored class relationships to fix node-specific allocation before prompt optimization.
\item A hidden-layer implementation with $C\times128$ trainable prompt parameters and no learned routing network.
\item Matched adaptation comparisons and shared-prompt ablations, with seed variability and paired uncertainty reported separately for each evaluated condition.
\end{itemize}
The contribution is the use of prototype relationships to organize hidden-layer prompts; it does not treat prototypes or node-dependent prompting as new primitives.

\section{Related Work}
\paragraph{Graph representations and pretraining.}
Graph convolution, neighborhood sampling, and graph attention provide complementary ways to construct node representations~\citep{gcn2017,sage2017,gat2018}. Self-supervised objectives include local--global information maximization~\citep{dgi2019}, augmentation-based contrastive learning~\citep{graphcl2020}, and masked reconstruction~\citep{graphmae2022}. Pretraining objectives can affect downstream transfer, and combining node- and graph-level objectives has been studied explicitly~\citep{pretrain2020}. Our study fixes the pretrained checkpoint on each graph and examines downstream adaptation, rather than proposing a new pretraining objective.

\paragraph{Aligning pretraining with downstream tasks.}
GPPT and GraphPrompt connect pretraining to downstream prediction through task reformulation and prompting~\citep{gppt2022,graphprompt2023}. All in One addresses multiple task levels with prompt tokens, structures, and insertion rules~\citep{allinone2023}; MultiGPrompt combines multiple pretext tasks and downstream prompts~\citep{multigprompt2024}. HGPrompt and HetGPT extend prompting to heterogeneous graph settings~\citep{hgprompt2024,hetgpt2024}. These methods address task or graph-type alignment. We focus on a narrower design question within few-shot node classification: how task-class references should allocate prompts across nodes.

\paragraph{Adaptive prompt allocation.}
GPF-plus learns input-dependent prompt combinations~\citep{gpf2023}, SUPT assigns prompts at subgraph level~\citep{supt2026}, and RELIEF learns node selection together with prompt content~\citep{relief2025}. ProNoG studies conditional node-specific prompts for non-homophilic graphs~\citep{pronog2025}, while DAGPrompT addresses distribution-aware adaptation~\citep{dagprompt2025}. Attribute-guided dynamic prompting has also been studied for robustness to graph perturbations~\citep{adprompt2026}. In contrast to learning an allocation function, \method{} computes support-conditioned prototype relationships once and keeps the resulting route fixed. These works delimit the design space; they are not all experimental baselines in our controlled comparison.

\paragraph{Class prototypes and evaluation.}
Prototypical Networks use support-set class representatives for prediction~\citep{proto2017}, and semi-supervised extensions incorporate unlabeled examples into prototype construction~\citep{semiproto2018}. We use these established ideas to obtain task references for routing. Routing prototypes live in aggregated attribute space; prediction prototypes come from the adapted encoder output. ProG~\citep{prog2024} motivates explicit evaluation protocols. \citet{mtg2026} report adaptation results across pretraining strategies; we reuse their baseline accuracy values as literature context while retaining local matched controls for the routing claim.

\section{Problem Setting}
Let $G=(V,E,X)$ be an attributed graph with $n$ nodes, features $X\in\mathbb R^{n\times d}$, and $C$ task classes. For each class $c$, a support set $S_c$ contains $k$ labeled nodes, with $k\in\{1,3,5\}$. Let $S=\bigcup_c S_c$ and let $T$ be a disjoint query set. The full graph topology and all node features are visible; only labels in $S$ are provided to fitting and hyperparameter selection. Query labels are used for scoring after predictions are sealed.

The encoder is a pretrained two-layer GCN~\citep{gcn2017} with hidden and output dimensions of 128. Its parameters are frozen for \method. The task is transductive, same-graph classification over known classes: query features participate in graph aggregation, prototype refinement, and global centering. It is not an inductive unseen-class evaluation. Throughout the paper, a \emph{condition} fixes the dataset, shot count, method, and pretrained checkpoint.

\section{Prototype-Routed Prompting}
\label{sec:method}
\begin{figure*}[t]
\centering
\fbox{\begin{minipage}[c][32mm][c]{0.92\linewidth}
\centering
\lang{Architecture figure placeholder}{主架构图占位符}\\[3mm]
\lang{Graph attributes and support labels $\longrightarrow$ anchored prototypes $\longrightarrow$ fixed routing}{图属性与支持标签 $\longrightarrow$ 支持集锚定原型 $\longrightarrow$ 固定路由}\\[3mm]
\lang{Fixed routing + trainable prompt tokens $\longrightarrow$ frozen GCN $\longrightarrow$ prototype prediction}{固定路由与可学习提示 $\longrightarrow$ 冻结 GCN $\longrightarrow$ 原型分类}
\end{minipage}}
\caption{\lang{Overview of \method. Support labels initialize class references in aggregated attribute space. Unlabeled graph features refine these references, and the resulting node--prototype associations determine a fixed routing matrix. Only the class prompt tokens are optimized; the routed prompts enter the first hidden layer of a frozen GCN. Query labels are excluded from fitting and selection. The final architecture artwork is intentionally left as a placeholder.}{\method 总体流程。支持标签在聚合属性空间中初始化类别参照，未标注图特征用于校正参照，节点与原型的关联决定固定路由矩阵。仅类别提示向量参与优化，路由后的提示注入冻结 GCN 的第一隐藏层。拟合与选参不使用查询标签。主架构图按要求暂用占位符。}}
\label{fig:architecture}
\end{figure*}

\paragraph{Overview.}
The method separates allocation from content. First, support-anchored prototypes define class relationships in aggregated attribute space. Second, a fixed routing matrix turns those relationships into node-specific combinations of prompts. Third, support supervision learns the prompt tokens inside the frozen encoder. The following subsections follow this computation order.
\subsection{Support-anchored class references}
Let $\widehat A=\widetilde D^{-1/2}(A+I)\widetilde D^{-1/2}$ denote the normalized adjacency with self-loops. We form a routing representation from two-hop aggregated attributes:
\begin{equation}
\begin{aligned}
X^{(2)}&=\widehat A^2X,\\
R&=\rownorm\!\left(X^{(2)}-\mathbf1\,\mean_{v\in V}X^{(2)}_v\right).
\end{aligned}
\label{eq:attributes}
\end{equation}
Here $\rownorm$ applies Euclidean normalization to each row. Routing therefore uses graph attributes and structure, rather than a separate learned network. Initial class references are
\begin{equation}
\mu_c^{(0)}=\frac{1}{k}\sum_{i\in S_c}R_i.
\label{eq:anchor}
\end{equation}

To incorporate the unlabeled graph while retaining a labeled reference, let $U_V=V\setminus S$. Starting from $\mu^{(0)}$, we use $i\in U_V$ and $t=0,1,2$ for three fixed refinement iterations:
\begin{align}
a_{ic}^{(t)}&=\frac{\exp(R_i^\top\norm{\mu_c^{(t)}}/0.05)}{\sum_{c'}\exp(R_i^\top\norm{\mu_{c'}^{(t)}}/0.05)},\label{eq:assignment}\\
\mu_c^{(t+1)}&=0.5\mu_c^{(0)}+0.5\frac{\sum_{i\in U_V}a_{ic}^{(t)}R_i}{\max(\sum_{i\in U_V}a_{ic}^{(t)},10^{-8})}.\label{eq:refinement}
\end{align}
The anchor is always the original support prototype, not the previous iterate. This is a fixed prototype construction; we do not claim an expectation-maximization objective or a convergence guarantee. No labels from $U_V$ enter these computations.

\subsection{From class relationships to a fixed routing matrix}
Using the refined prototypes, we compute node--class associations at temperature $\tau_r=0.1$:
\begin{equation}
\begin{aligned}
Q_{ic}&=\softmax_c\!\left(R_i^\top\norm{\mu_c^{(3)}}/\tau_r\right),\\
B&=\widehat A\left(Q-\mathbf1\,\mean_{v\in V}Q_v\right).
\end{aligned}
\label{eq:routing}
\end{equation}
Centering expresses a node's class association relative to the graph-wide average; neighborhood aggregation incorporates adjacent association patterns. $Q$ is a row-normalized association matrix, whereas $B$ contains signed prompt-combination coefficients. We do not interpret $B$ as a class-probability matrix. $B$ is computed once per support set and kept fixed throughout training.

\subsection{Adapting frozen hidden computations}
Let $H^{\mathrm{pre}}$ be the first GCN layer's preactivation on the original graph. We learn a zero-initialized matrix $P\in\mathbb R^{C\times128}$. With $s$ the vector of channelwise sample standard deviations of $H^{\mathrm{pre}}$, clamped below by $10^{-6}$, the prompted hidden representation is
\begin{equation}
\widetilde H^{\mathrm{pre}}=H^{\mathrm{pre}}+10\,[B\tanh(P)]\odot s.
\label{eq:injection}
\end{equation}
The scale vector is broadcast over nodes. The existing LeakyReLU and second graph convolution then produce $Z$. Thus, prompt content is learned through support supervision, while routing and backbone parameters remain fixed. Zero initialization recovers the frozen encoder at the start of optimization. The trainable prompt counts are 896, 768, and 384 for Cora, CiteSeer, and PubMed.

\subsection{Training and prediction}
We optimize the prompt matrix for 100 Adam steps with learning rate 0.01 and no weight decay:
\begin{equation}
\begin{aligned}
\mathcal L(P)={}&\frac{1}{|S|}\sum_{i\in S}\operatorname{CE}(\ell_i(P),y_i)\\
&+10^{-3}\mean(P\odot P).
\end{aligned}
\label{eq:loss}
\end{equation}
The support logits $\ell_i$ use cosine prototype classification. The same prototype readout is kept throughout the routing ablation. Support leave-one-out construction, prediction normalization, and temperatures are specified in Appendix~\ref{app:protocol}.

The routing is a one-time computation. For sparse graphs, the two attribute aggregation steps cost $O(|E|d)$; prototype refinement costs $O(nCd)$ per iteration; and multiplying $B$ by prompt tokens costs $O(nC\cdot128)$ per optimization step, in addition to the frozen encoder's forward and backward computations. Freezing weights reduces trainable parameters but does not eliminate gradient computation through the second convolution. We report parameter counts, not an unmeasured runtime advantage.

\section{Experiments}
\subsection{Protocol and comparison methods}
\label{sec:protocol}
\begin{table*}[t]
\centering
\linespread{1}\selectfont
\captionsetup{font=small,skip=5pt}
\caption{\lang{\textbf{Accuracy on few-shot node classification.} Higher is better ($\uparrow$).}{\textbf{少样本节点分类的 Accuracy。}越高越好（$\uparrow$）。}}
\label{tab:main-accuracy}
\fontsize{9}{11}\selectfont
\setlength{\tabcolsep}{1pt}
\renewcommand{\arraystretch}{1.18}
\setlength{\heavyrulewidth}{0.8pt}
\setlength{\lightrulewidth}{0.4pt}
\begin{tabular*}{\linewidth}{@{\extracolsep{\fill}}l*{9}{r}@{}}
\toprule
\lang{Method}{方法} & \multicolumn{3}{c}{\textbf{Cora}} & \multicolumn{3}{c}{\textbf{CiteSeer}} & \multicolumn{3}{c}{\textbf{PubMed}} \\
\cmidrule(lr){2-4}\cmidrule(lr){5-7}\cmidrule(l){8-10}
 & 1-shot & 3-shot & 5-shot & 1-shot & 3-shot & 5-shot & 1-shot & 3-shot & 5-shot \\
\midrule
\multicolumn{10}{@{}l}{\textit{\lang{Literature-reported baselines}{文献报告基线}~\citep{mtg2026}}} \\[2pt]
Supervised & \meanstd{26.56}{5.55} & \meanstd{37.79}{9.16} & \meanstd{50.25}{8.37} & \meanstd{21.78}{7.32} & \meanstd{35.18}{6.86} & \meanstd{41.22}{6.30} & \meanstd{39.37}{16.34} & \meanstd{57.33}{4.64} & \meanstd{67.88}{2.18} \\
Fine-tuning & \meanstd{52.61}{1.73} & \meanstd{51.97}{2.84} & \meanstd{62.66}{3.55} & \meanstd{35.05}{4.37} & \meanstd{45.08}{2.09} & \meanstd{39.54}{3.54} & \meanstd{46.74}{14.89} & \meanstd{65.40}{3.00} & \bestmeanstd{70.91}{4.87} \\
\midrule
GPPT & \meanstd{43.15}{9.44} & \meanstd{43.84}{6.11} & \meanstd{51.98}{3.43} & \meanstd{37.26}{6.17} & \meanstd{42.34}{8.31} & \meanstd{45.77}{7.41} & \secondmeanstd{48.31}{17.72} & \meanstd{67.43}{2.96} & \meanstd{66.97}{3.70} \\
Gprompt & \bestmeanstd{56.66}{11.22} & \bestmeanstd{63.78}{5.77} & \bestmeanstd{69.03}{3.61} & \secondmeanstd{53.21}{10.94} & \secondmeanstd{60.00}{6.18} & \secondmeanstd{66.13}{1.64} & \meanstd{39.74}{15.35} & \meanstd{66.68}{3.53} & \meanstd{67.87}{2.08} \\
All-in-one & \meanstd{52.39}{10.17} & \meanstd{48.09}{4.83} & \meanstd{30.36}{13.48} & \meanstd{40.41}{2.80} & \meanstd{48.09}{8.18} & \meanstd{27.93}{10.59} & \meanstd{45.17}{6.45} & \meanstd{65.79}{5.79} & \meanstd{46.16}{15.83} \\
GPF & \meanstd{38.57}{5.41} & \meanstd{34.84}{19.83} & \meanstd{35.43}{1.02} & \meanstd{31.16}{8.05} & \meanstd{25.92}{12.30} & \meanstd{25.12}{3.01} & \bestmeanstd{49.99}{8.86} & \bestmeanstd{71.20}{2.82} & \meanstd{68.96}{3.99} \\
GPF-plus & \secondmeanstd{55.77}{10.30} & \secondmeanstd{56.38}{5.37} & \secondmeanstd{66.22}{6.20} & \bestmeanstd{59.67}{11.87} & \bestmeanstd{72.48}{5.63} & \bestmeanstd{75.73}{2.19} & \meanstd{46.64}{18.97} & \secondmeanstd{70.85}{4.03} & \secondmeanstd{69.59}{4.33} \\
\midrule
\multicolumn{10}{@{}l}{\textit{\lang{Local measurements (different protocol)}{本地测量（协议不同）}}} \\[2pt]
\frozenname & \meanstd{23.17}{5.48} & \meanstd{29.42}{3.47} & \meanstd{33.56}{5.03} & \meanstd{21.14}{2.87} & \meanstd{25.29}{2.16} & \meanstd{27.38}{2.84} & \secondmeanstd{55.44}{11.42} & \secondmeanstd{65.55}{8.11} & \secondmeanstd{68.37}{3.41} \\
\ftname & \secondmeanstd{45.41}{10.18} & \secondmeanstd{61.54}{6.78} & \secondmeanstd{67.50}{5.06} & \secondmeanstd{35.76}{6.80} & \secondmeanstd{42.22}{5.88} & \secondmeanstd{47.89}{4.85} & \meanstd{48.77}{8.54} & \meanstd{59.75}{7.97} & \meanstd{63.96}{5.18} \\
\midrule\rowcolor{black!5}
\textbf{\method} & \bestmeanstd{54.59}{9.18} & \bestmeanstd{70.76}{2.76} & \bestmeanstd{73.02}{1.95} & \bestmeanstd{47.12}{6.31} & \bestmeanstd{61.20}{4.98} & \bestmeanstd{66.17}{2.16} & \bestmeanstd{58.70}{9.78} & \bestmeanstd{68.66}{4.58} & \bestmeanstd{70.98}{2.46} \\
\bottomrule
\end{tabular*}
\par\vspace{3pt}
\begin{minipage}{\linewidth}\fontsize{8}{9.5}\selectfont\raggedright \lang{Mean $\pm$ std. (\%). Literature panel: \citet{mtg2026}, Tables 4--6, five trials, 90\% test allocation, best reported across pretraining strategies. Local panel: fixed edge-pretrained checkpoint and official test split; three training seeds averaged per support draw, then mean/sample std. over 20 draws. \textbf{Bold}/\underline{underline}: best/second within each panel and column only. Protocols differ; no cross-panel ranking, pairing, or pooled average.}{均值 $\pm$ 标准差（\%）。文献面板取自 \citet{mtg2026} 表 4--6：5 次抽样、90\% 测试集、不同预训练策略中的最佳报告结果。本地面板使用固定边预测预训练权重与官方测试划分，先平均每个支持抽样的 3 个训练 seed，再统计 20 次抽样的均值／样本标准差。\textbf{粗体}／\underline{下划线}仅表示同面板同列最优／次优；不跨协议排名、配对检验或求平均。}\end{minipage}
\end{table*}

\begin{table*}[t]
\centering
\linespread{1}\selectfont
\captionsetup{font=small,skip=5pt}
\caption{\lang{\textbf{Macro-F1 on few-shot node classification.} Higher is better ($\uparrow$).}{\textbf{少样本节点分类的 Macro-F1。}越高越好（$\uparrow$）。}}
\label{tab:main-macro_f1}
\fontsize{9}{11}\selectfont
\setlength{\tabcolsep}{1pt}
\renewcommand{\arraystretch}{1.18}
\setlength{\heavyrulewidth}{0.8pt}
\setlength{\lightrulewidth}{0.4pt}
\begin{tabular*}{\linewidth}{@{\extracolsep{\fill}}l*{9}{r}@{}}
\toprule
\lang{Method}{方法} & \multicolumn{3}{c}{\textbf{Cora}} & \multicolumn{3}{c}{\textbf{CiteSeer}} & \multicolumn{3}{c}{\textbf{PubMed}} \\
\cmidrule(lr){2-4}\cmidrule(lr){5-7}\cmidrule(l){8-10}
 & 1-shot & 3-shot & 5-shot & 1-shot & 3-shot & 5-shot & 1-shot & 3-shot & 5-shot \\
\midrule
\frozenname & \meanstd{22.51}{4.82} & \meanstd{29.07}{3.47} & \meanstd{33.07}{4.47} & \meanstd{20.39}{2.62} & \meanstd{24.32}{1.87} & \meanstd{26.45}{2.71} & \secondmeanstd{53.44}{10.49} & \secondmeanstd{63.19}{8.60} & \secondmeanstd{66.27}{3.55} \\
\ftname & \secondmeanstd{44.38}{10.07} & \secondmeanstd{60.94}{6.41} & \secondmeanstd{66.99}{4.52} & \secondmeanstd{30.99}{6.30} & \secondmeanstd{38.43}{5.80} & \secondmeanstd{44.62}{4.98} & \meanstd{44.52}{9.85} & \meanstd{57.64}{9.26} & \meanstd{62.53}{6.05} \\
\midrule
GPF & \meanstd{7.08}{2.07} & \meanstd{21.14}{6.64} & \meanstd{25.87}{6.19} & \meanstd{6.78}{1.56} & \meanstd{13.81}{5.45} & \meanstd{15.27}{5.85} & \meanstd{25.62}{6.56} & \meanstd{46.88}{9.09} & \meanstd{52.42}{6.87} \\
GPF-plus & \meanstd{27.08}{6.92} & \meanstd{39.31}{4.55} & \meanstd{51.82}{5.06} & \meanstd{12.83}{4.10} & \meanstd{26.16}{2.85} & \meanstd{29.46}{2.64} & \meanstd{34.89}{8.06} & \meanstd{54.47}{4.51} & \meanstd{58.79}{5.19} \\
\gpromptname & \meanstd{21.03}{4.36} & \meanstd{26.98}{3.41} & \meanstd{30.82}{3.43} & \meanstd{18.00}{2.95} & \meanstd{23.71}{2.39} & \meanstd{26.72}{2.49} & \meanstd{45.35}{10.95} & \meanstd{59.93}{5.76} & \meanstd{64.79}{4.02} \\
\midrule\rowcolor{black!5}
\textbf{\method} & \bestmeanstd{53.34}{8.94} & \bestmeanstd{69.38}{2.91} & \bestmeanstd{72.04}{1.77} & \bestmeanstd{44.26}{5.75} & \bestmeanstd{58.19}{4.67} & \bestmeanstd{63.03}{2.12} & \bestmeanstd{57.37}{10.07} & \bestmeanstd{68.21}{4.08} & \bestmeanstd{70.39}{2.47} \\
\bottomrule
\end{tabular*}
\par\vspace{3pt}
\begin{minipage}{\linewidth}\fontsize{8}{9.5}\selectfont\raggedright \lang{Statistics and emphasis follow Table~\ref{tab:main-accuracy}. Each column is a separate dataset/shot condition.}{统计与强调规则同表~\ref{tab:main-accuracy}；每列对应独立的数据集／shot 条件。}\lang{ $\dagger$: GraphPrompt under the common-protocol adaptation.}{ $\dagger$：GraphPrompt 的统一协议适配实现。}\end{minipage}
\end{table*}

\begin{table*}[t]
\centering
\linespread{1}\selectfont
\captionsetup{font=small,skip=5pt}
\caption{\lang{\textbf{Prototype-routing ablation.} Matched prompt recipe and readout.}{\textbf{原型路由消融。}提示训练配方与读出保持一致。}}
\label{tab:ablation}
\fontsize{9}{11}\selectfont
\setlength{\tabcolsep}{1pt}
\renewcommand{\arraystretch}{1.18}
\setlength{\heavyrulewidth}{0.8pt}
\setlength{\lightrulewidth}{0.4pt}
\begin{tabular*}{\linewidth}{@{\extracolsep{\fill}}l*{9}{r}@{}}
\toprule
\lang{Method}{方法} & \multicolumn{3}{c}{\textbf{Cora}} & \multicolumn{3}{c}{\textbf{CiteSeer}} & \multicolumn{3}{c}{\textbf{PubMed}} \\
\cmidrule(lr){2-4}\cmidrule(lr){5-7}\cmidrule(l){8-10}
 & 1-shot & 3-shot & 5-shot & 1-shot & 3-shot & 5-shot & 1-shot & 3-shot & 5-shot \\
\midrule
\multicolumn{10}{@{}l}{\textit{Accuracy (\%) $\uparrow$}} \\[2pt]
\frozenname & \secondmeanstd{24.64}{4.63} & \meanstd{30.86}{3.27} & \meanstd{35.16}{2.40} & \secondmeanstd{21.69}{3.47} & \secondmeanstd{24.96}{3.51} & \secondmeanstd{27.52}{3.72} & \secondmeanstd{53.72}{10.60} & \secondmeanstd{62.49}{9.18} & \secondmeanstd{65.68}{5.56} \\
\sharedname & \meanstd{24.15}{4.41} & \secondmeanstd{31.25}{3.68} & \secondmeanstd{36.30}{2.72} & \meanstd{20.64}{3.07} & \meanstd{23.71}{3.88} & \meanstd{25.60}{3.82} & \meanstd{52.61}{10.70} & \meanstd{60.89}{9.03} & \meanstd{65.00}{5.72} \\
\midrule\rowcolor{black!5}
\textbf{\method} & \bestmeanstd{53.80}{8.45} & \bestmeanstd{69.97}{3.53} & \bestmeanstd{72.87}{3.36} & \bestmeanstd{43.92}{11.01} & \bestmeanstd{59.67}{5.98} & \bestmeanstd{64.70}{3.53} & \bestmeanstd{59.77}{7.47} & \bestmeanstd{65.09}{8.21} & \bestmeanstd{68.34}{6.97} \\
\midrule
\multicolumn{10}{@{}l}{\textit{Macro-F1 (\%) $\uparrow$}} \\[2pt]
\frozenname & \secondmeanstd{23.79}{4.61} & \secondmeanstd{30.11}{2.90} & \meanstd{34.48}{2.30} & \secondmeanstd{20.84}{3.15} & \secondmeanstd{23.85}{3.32} & \secondmeanstd{26.11}{3.34} & \secondmeanstd{52.29}{9.78} & \secondmeanstd{61.43}{9.05} & \secondmeanstd{64.19}{5.48} \\
\sharedname & \meanstd{22.78}{4.10} & \meanstd{30.03}{3.70} & \secondmeanstd{34.69}{2.61} & \meanstd{19.75}{2.74} & \meanstd{22.28}{3.34} & \meanstd{24.42}{3.38} & \meanstd{51.98}{10.37} & \meanstd{59.21}{10.06} & \meanstd{63.59}{6.09} \\
\midrule\rowcolor{black!5}
\textbf{\method} & \bestmeanstd{51.82}{7.47} & \bestmeanstd{68.37}{4.08} & \bestmeanstd{72.02}{2.74} & \bestmeanstd{41.48}{10.31} & \bestmeanstd{56.34}{5.54} & \bestmeanstd{61.34}{3.09} & \bestmeanstd{59.39}{7.93} & \bestmeanstd{64.82}{8.80} & \bestmeanstd{67.80}{7.59} \\
\bottomrule
\end{tabular*}
\par\vspace{3pt}
\begin{minipage}{\linewidth}\fontsize{8}{9.5}\selectfont\raggedright \lang{Local original-ProG runs with support/test overlap: mean $\pm$ sample std. over 20 tasks after averaging three training seeds per task. \textbf{Bold}/\underline{underline}: best/second in each condition. No cross-condition average.}{本地原ProG实现，保留支持／测试重叠：先平均每任务3个训练seed，再报告20个任务的均值 $\pm$ 样本标准差。\textbf{粗体}／\underline{下划线}表示同条件最优／次优；不跨条件平均。}\lang{ Frozen: no prompt; shared prompt: uniform allocation; PRPrompt: prototype routing.}{ 冻结预训练：无提示；共享提示：均匀分配；PRPrompt：原型路由。}\end{minipage}
\end{table*}

Our local evaluation covers Cora, CiteSeer, and PubMed at 1, 3, and 5 shots per class. Official training masks define support pools and official test masks define fixed query sets. A shared per-class permutation yields nested support sets. Each local condition contains 20 support seeds and three training seeds, with matched graph, checkpoint, and support/query identities. The pretrained source is the original edge-prediction checkpoint for each dataset; its pretraining seed is unavailable.

The local references are \textbf{Frozen}, with no prompt and the same final prototype readout as \method, and \textbf{Full FT}, updating the pretrained encoder and a task readout. Prompt baseline accuracy is taken from Tables 4--6 of \citet{mtg2026}: GPPT, Gprompt (GraphPrompt), All-in-one, GPF, and GPF-plus, together with their supervised and fine-tuning references. Their MTG method is excluded. The source combines ProG results with author reproductions and reports best performance across pretraining strategies; the selected strategy is not identified for each cell.

\paragraph{Literature protocol.}
The source uses a 90\% test allocation, five independent task samplings, and random hyperparameter search. Its main backbone is a two-layer GCN with hidden dimension 128. These results are not matched to our test split, checkpoint, or fitting budget. Table~\ref{tab:main-accuracy} therefore separates literature and local panels. We retain the source means and standard deviations without reaggregation, and compute no paired intervals or global ranking across panels. The source tables report Accuracy only; the other metric tables contain local measurements.

\method{} and the shared-prompt ablation use one fixed recipe. Local Full FT uses three support-only folds at 3 and 5 shots and four candidates: learning rate in $\{0.001,0.01\}$ and weight decay in $\{0,0.0001\}$. Each candidate fit and final fit uses 100 steps; 1-shot uses fixed parameters. Consequently, local methods do not have identical total optimization budgets. Appendix~\ref{app:protocol} gives the remaining settings.

\paragraph{Statistical unit.}
For a fixed local condition, we first average three training seeds within each support draw, then report the mean and sample standard deviation over the 20 support draws. We never average performance across datasets or shot counts. Local paired intervals resample support draws with 10,000 bootstrap samples. They are unadjusted 95\% intervals for the fixed graph and checkpoint, and apply only to Frozen, Full FT, and shared-prompt comparisons. Accuracy and Macro-F1 are the primary reported metrics; the remaining local metrics appear in Appendix~\ref{app:more}.

\subsection{RQ1: Adaptation performance}
The local comparison tests adaptation under the same graph, checkpoint, and labeled support budget. Table~\ref{tab:main-accuracy} adds literature context with a separate protocol; Table~\ref{tab:main-macro_f1} reports local Macro-F1. Each column is an individual condition, without a pooled average.

On Cora at 5 shots, \method{} reaches 73.02\% accuracy compared with 33.56\% for Frozen and 67.50\% for Full FT. On CiteSeer at the same shot count, the corresponding accuracies are 66.17\%, 27.38\%, and 47.89\%. PubMed has a stronger frozen reference: 68.37\% at 5 shots, compared with 70.98\% for \method{} and 63.96\% for Full FT. Thus, the improvement over frozen pretraining is substantially smaller on PubMed, an important dependence that a cross-dataset average would hide.

\paragraph{Literature-reported prompt results.}
The imported values do not support a claim of uniformly higher accuracy. On CiteSeer, GPF-plus reports 59.67, 72.48, and 75.73\% at 1, 3, and 5 shots, whereas our corresponding measurements are 47.12, 61.20, and 66.17\%. On Cora at 5 shots, our 73.02\% lies above the reported Gprompt value of 69.03\%; on PubMed at 3 shots, our 68.66\% lies below the reported GPF value of 71.20\%. These are descriptive cross-source comparisons, not matched estimates of method effects. The evidence for routing comes from the local ablation below; local paired differences are provided in Appendix~\ref{app:paired}.

\subsection{RQ2: Value of prototype routing}
The central ablation replaces $B$ in Eq.~\eqref{eq:injection} by the constant matrix with entries $1/C$. This retains the prompt matrix, scale, injection point, optimization, and final readout, but removes node-dependent allocation. The constant matrix is used directly; centering it would erase the prompt. We also retain the no-prompt Frozen reference.

Table~\ref{tab:ablation} shows that a shared prompt alone does not account for the gains. At 5 shots, routing improves accuracy over the shared prompt by 36.94 percentage points on Cora, 40.45 on CiteSeer, and 3.02 on PubMed. The respective paired 95\% intervals are $[34.87,39.06]$, $[39.06,41.84]$, and $[1.39,4.51]$ percentage points. The improvement is also positive in both primary metrics in each 1- and 3-shot condition.

This comparison supports the complete prototype-routing construction relative to a shared prompt. It does not isolate prototype refinement, centering, and graph aggregation as separate causes. A uniform mixture also reduces the effective degrees of freedom despite retaining the nominal prompt parameter count. We therefore do not interpret the comparison as a capacity-matched proof of a unique routing mechanism.

\section{Limitations}
The evaluation covers three citation graphs, one fixed pretrained checkpoint per graph, and same-graph known-class classification. Pretraining seeds are unknown, so variation across pretraining runs is not estimated. The reporting scope was chosen after the experiments were completed; the analysis is descriptive rather than a new independent confirmatory study. It cannot support claims about graph classification, unseen classes, or arbitrary pretraining sources.

Literature baselines use a different split, five trials, and selection across pretraining strategies. Their aggregate values provide context but cannot isolate an adaptation-method effect against our fixed-checkpoint runs or support paired significance tests. Local fine-tuning selection also incurs a different budget from the fixed \method{} recipe. The routing ablation changes effective prompt capacity and does not distinguish every component of the routing construction. Finally, the method uses query features and graph-wide statistics; deployment without such transductive access requires separate evaluation.

\section{Conclusion}
\method{} uses support-anchored prototypes to organize the allocation of trainable hidden-layer prompts in a frozen graph encoder. In the evaluated Cora, CiteSeer, and PubMed settings, condition-wise comparisons and shared-prompt ablations support the value of this allocation. The evidence favors prototype routing as a practical adaptation design within the tested protocol, while leaving its generality across graph families and pretraining sources open.

\section*{Broader Impact}
The study concerns public citation-network benchmarks and does not evaluate decisions about individuals. Errors can propagate through graph aggregation and unlabeled prototype refinement, particularly when the graph or class structure differs from the evaluated setting. Applications involving consequential decisions require domain-specific evaluation beyond these benchmarks.
